 \documentclass[a4paper,10pt]{article}
\include{amsymb}
\include{marvosym}

\usepackage[francais]{babel}
\usepackage[cyr]{aeguill}
\usepackage[applemac]{inputenc}
\usepackage{graphicx}
\usepackage{xspace}
\usepackage[a4paper]{geometry}
\usepackage{latexsym,amsmath,amssymb,textcomp}
\usepackage{moreverb}
\usepackage{listings}
\usepackage{multirow}
\usepackage{titling}
\usepackage{mathabx}

\usepackage{pdfpages}

\usepackage{hyperref}


\newlength{\indentationnota}
\newlength{\largeurlignenota}
\newlength{\paddingnota}
\newlength{\largeurnota}
\setlength{\paddingnota}{5pt}
\setlength{\largeurnota}{0.9cm}

\definecolor{pink}{rgb}{1,0.5,1}
\makeatletter

\newenvironment{pictonote}[1]{% on passe le nom du �chier en argument
\begin{list}{}{%
\setlength{\labelsep}{0pt}%
\setlength{\rightmargin}{15pt}%
\setlength{\paddingnota}{5pt}}
\item%
\setlength{\indentationnota}{%
\@totalleftmargin+\largeurnota+\paddingnota}%
\setlength{\largeurlignenota}{%
\linewidth-\largeurnota-\paddingnota}%
\parshape=3%
\indentationnota\largeurlignenota%
\indentationnota\largeurlignenota%
\@totalleftmargin\linewidth%
\raisebox{-\largeurnota+2.2ex}[0pt][0pt]{%
\makebox[0pt][r]{%
\includegraphics[width=\largeurnota]{#1}%
\hspace{\paddingnota}}}%
\ignorespaces}{%
\end{list}}

\makeatother

\newenvironment{attention}%
{\begin{pictonote}{/Users/benjamin/Documents/Education/LaTeX/danger}}%
{\end{pictonote}}

\newenvironment{unicorn}%
{\begin{pictonote}{/Users/benjamin/Documents/Education/LaTeX/angry_unicorn}}%
{\end{pictonote}}

\newcommand\unicornbox[1]{\colorbox{pink}{\color{white}{#1}}}

\newcommand\cf{\emph{c.f}}

\newcommand\signature{%
\begin{figure}[br]
	\begin{flushright}
	\begin{minipage}{8cm}
		\begin{center}
		\reflectbox{\includegraphics[width=8cm]{/Users/benjamin/Documents/Education/LaTeX/unicorn}}\\
		\theauthor
		\end{center}
	\end{minipage}
	\end{flushright}
\end{figure}}

\newcommand\danger{\raisebox{-0.4ex}{\LARGE $\triangle$ \normalsize} \hspace{-4.1ex}! \hspace{1ex}}
\newcommand\PL{Programmation Logique\xspace}
\newcommand\pl{\bsc{Prolog}\xspace}


\title{Rdf - TP5 : Classement de caract�re manuscrits}
\author{Fran�ois \bsc{Lepan} - Benjamin \bsc{Van Ryseghem}}

\begin{document}
\maketitle 

\section{Visualiser les Donnees}

\begin{verbatimtab}
y = matrix(train0(1,:),28,28)';
Matplot(y);
\end{verbatimtab}

Afin de de visualiser les autres chiffres manuscrits il suffit de changer le nom de train0 en train$N$ o� $N$ est le chiffre � visualiser.

\section{L'algorithme de plus proche voisin ou kNN}

\subsection{Est-ce que les r�sultats changent si $\displaystyle d^2(x, y)=\Sigma^n_{i=1}(xi - yi)^2$ est utilis�e comme distance?}

Non car tout ce que l'on fait c'est comparer ces distances ($\geq 0$) donc leurs valeurs �lev� au carr� ne changeront rien lors de la comparaison.

Ex: si a = 2 et b = 3

$a < b \Rightarrow a^2 < b^2$

$2 < 3 \Rightarrow 4 < 9$

\subsection{Quelles sont les erreurs moyennes de classification pour k = 1, 3, 5, 10?}

Pour k = 1 on a une erreur moyenne de 0.6885923. 

�tant donn� que la plupart des gens on 0.1 je me dit qu'il doit y avoir une erreur dans notre code mais je n'arrive pas � savoir o�.

\subsection{Cette approche est-elle supervis�e ?}

Oui cette approche est supervis�e, car il faut l'alimenter de donn�es d'entrainement afin de savoir comment classifier le nouvel �l�ment entrant.

\section{L�algorithme Na�ve Bayes}

\subsection{Quelle est l�hypoth�se principale dans cette approche?}

$P(X = (x1,�,xm~|~Ci)) = \prod_{j=1}^m P(xj | Ci)$ \\
\\
 Propri�t�: L'ind�pendance conditionnelle des attributs donc on peut faire le produit de ses �l�ments.
 
\subsection{Calcul des erreurs de classification}

???????????????????

\subsection{Cette approche est-elle supervis�e ?}

Oui cette approche est supervis�e, car elle utilise des donn�es d'entrainement tout comme kNN.

\subsection{Cette approche est-elle probabiliste ?}

Oui car elle est bas�e sur la fr�quence de chaque valeur (ici 0 $\rightarrow$ 255) quelque soit le point d�entrainement et quel que soit l�emplacement de cette valeur dans le point.

\end{document}